% ============================================================
% Resumen de Trabajo Final
% Archivo autónomo para Overleaf. Compilador: pdfLaTeX.
% Copiar todo este contenido dentro de main.tex.
% ============================================================
\documentclass[11pt,a4paper]{article}

% Idioma, codificación y tipografía
\usepackage[utf8]{inputenc}
\usepackage[T1]{fontenc}
\usepackage[spanish,es-nodecimaldot]{babel}
\usepackage{lmodern}
\usepackage{microtype}
\DeclareUnicodeCharacter{2192}{\ensuremath{\rightarrow}}
\DeclareUnicodeCharacter{2193}{\ensuremath{\downarrow}}

% Página y estructura
\usepackage[a4paper,margin=2.4cm,headheight=15pt]{geometry}
\usepackage{parskip}
\usepackage{enumitem}
\usepackage{titlesec}
\usepackage{fancyhdr}

% Matemática
\usepackage{amsmath,amssymb}

% Colores, enlaces y bloques de código
\usepackage{xcolor}
\usepackage{hyperref}
\usepackage{xurl}
\usepackage{listings}

% Paleta visual
\definecolor{ink}{HTML}{17263A}
\definecolor{muted}{HTML}{52657A}
\definecolor{accent}{HTML}{167D9A}
\definecolor{soft}{HTML}{F1F6F8}
\definecolor{line}{HTML}{CBD6DF}

% Enlaces y metadatos
\hypersetup{
  colorlinks=true,
  linkcolor=accent,
  urlcolor=accent,
  citecolor=accent,
  pdftitle={Resumen de Trabajo Final},
  pdfsubject={Simulación progresiva de un circuito neuromotor},
  pdfauthor={Neural Network Exploration}
}
\urlstyle{same}

% Encabezados
\titleformat{\section}
  {\Large\bfseries\color{ink}}{}{0pt}{}[\vspace{1mm}\color{line}\titlerule]
\titleformat{\subsection}
  {\large\bfseries\color{accent}}{}{0pt}{}
\titlespacing*{\section}{0pt}{1.1em}{0.55em}
\titlespacing*{\subsection}{0pt}{0.9em}{0.35em}
\setcounter{secnumdepth}{0}
\setcounter{tocdepth}{2}

% Listas
\setlist[itemize]{leftmargin=1.6em,itemsep=0.25em,topsep=0.35em}
\setlist[enumerate]{leftmargin=1.8em,itemsep=0.35em,topsep=0.35em}

% Código y diagramas de texto. Las reglas literate permiten compilar
% directamente las flechas presentes en el Markdown con pdfLaTeX.
\lstdefinestyle{reportcode}{
  basicstyle=\ttfamily\small,
  backgroundcolor=\color{soft},
  frame=single,
  rulecolor=\color{line},
  framesep=6pt,
  xleftmargin=0.5em,
  xrightmargin=0.5em,
  breaklines=true,
  columns=fullflexible,
  keepspaces=true,
  showstringspaces=false,
  literate={↓}{{$\downarrow$}}1 {→}{{$\rightarrow$}}1
}
\lstset{style=reportcode}

% Compatibilidad con listas compactas producidas por Pandoc
\providecommand{\tightlist}{%
  \setlength{\itemsep}{0pt}\setlength{\parskip}{0pt}}

% Encabezado y pie
\pagestyle{fancy}
\fancyhf{}
\fancyhead[L]{\small\color{muted}Resumen de Trabajo Final}
\fancyhead[R]{\small\color{muted}Circuito neuromotor}
\fancyfoot[C]{\small\color{muted}\thepage}
\renewcommand{\headrulewidth}{0.4pt}
\renewcommand{\headrule}{\hbox to\headwidth{\color{line}\leaders\hrule height \headrulewidth\hfill}}

% Evita líneas desbordadas por URLs o nombres técnicos largos
\setlength{\emergencystretch}{3em}

\begin{document}

\begin{titlepage}
  \centering
  \vspace*{3.0cm}
  {\Huge\bfseries\color{ink} Resumen de Trabajo Final \par}
  \vspace{1.2cm}
  {\LARGE\bfseries\color{accent} Simulación progresiva de un circuito neuromotor mediante neuronas de Izhikevich \par}
  \vfill
  {\large Neural Network Exploration \par}
  \vspace{0.3cm}
  {\large \today \par}
\end{titlepage}

\tableofcontents
\newpage

\hypertarget{introducciuxf3n}{%
\section{1. Introducción}\label{introducciuxf3n}}

Este trabajo desarrolla una simulación computacional progresiva de un
circuito neuromotor simplificado. El recorrido comienza con una neurona
individual, continúa con pools neuronales y redes conectadas, incorpora
un circuito de inhibición recurrente tipo Renshaw y finaliza con la
transformación de los spikes de las motoneuronas en una señal muscular
estimada.

El propósito no es reproducir de manera completa la fisiología del
sistema motor ni construir un modelo biomecánico. Se busca una
implementación pequeña, reproducible, visual e interpretable que permita
comprender cómo se conectan distintos niveles de una simulación
neuromotora:

\begin{lstlisting}
Input externo
      ↓
Dinámica de una neurona spiking
      ↓
Actividad de un pool neuronal
      ↓
Conectividad sináptica
      ↓
Feedback inhibitorio tipo Renshaw
      ↓
Conversión de spikes en activación muscular simplificada
\end{lstlisting}

\hypertarget{objetivo-general}{%
\section{2. Objetivo general}\label{objetivo-general}}

Construir un modelo computacional progresivo que vincule la actividad de
neuronas spiking con un circuito de inhibición recurrente y una salida
muscular simplificada, priorizando claridad conceptual, modularidad y
reproducibilidad.

\hypertarget{objetivos-especuxedficos}{%
\section{3. Objetivos específicos}\label{objetivos-especuxedficos}}

\begin{itemize}
\item
  Implementar el modelo neuronal de Izhikevich
\item
  Explorar diferentes formas de corriente externa o comando motor
\item
  Simular una neurona individual y un pool neuronal heterogéneo
\item
  Representar conexiones sinápticas mediante matrices de pesos
\item
  Incorporar conexiones excitatorias e inhibitorias con dinámica
  temporal
\item
  Construir un circuito con motoneuronas y células de Renshaw
  idealizadas
\item
  Comparar configuraciones manteniendo controladas las fuentes de
  aleatoriedad
\item
  Convertir los spikes de las motoneuronas en una señal contráctil
  basada en twitches
\end{itemize}

\hypertarget{alcance-del-modelo}{%
\section{4. Alcance del modelo}\label{alcance-del-modelo}}

El trabajo se plantea como una maqueta computacional. Las neuronas,
sinapsis, conexiones y respuestas musculares utilizan unidades y
parámetros abstractos o idealizados. El circuito permite estudiar
relaciones funcionales entre sus componentes, pero no pretende estimar
directamente voltajes, corrientes o fuerzas de un sistema biológico
real.

Las principales simplificaciones metodológicas son:

\begin{itemize}
\tightlist
\item
  parámetros neuronales no calibrados con registros experimentales
  específicos
\item
  conexiones aleatorias
\item
  comando motor artificial
\item
  sinapsis basadas en corriente y estados exponenciales
\item
  ausencia de retardos axonales, plasticidad y potenciales de reversión
\item
  ausencia de feedback sensorial
\item
  ausencia de músculo antagonista y biomecánica articular
\item
  fuerza muscular expresada en unidades arbitrarias
\end{itemize}

\hypertarget{modelo-neuronal}{%
\section{5. Modelo neuronal}\label{modelo-neuronal}}

\hypertarget{ecuaciones-de-izhikevich}{%
\subsection{5.1. Ecuaciones de
Izhikevich}\label{ecuaciones-de-izhikevich}}

La dinámica neuronal se representa mediante:

\$ \frac{dv}{dt}=0.04v\^{}2+5v+140-u+I \$

\$ \frac{du}{dt}=a(bv-u) \$

donde:

\begin{itemize}
\tightlist
\item
  \texttt{v} es el potencial de membrana;
\item
  \texttt{u} es una variable de recuperación;
\item
  \texttt{I} es la corriente total recibida;
\item
  \texttt{a} controla la velocidad de recuperación;
\item
  \texttt{b} controla el acoplamiento entre \texttt{v} y \texttt{u};
\item
  \texttt{c} es el potencial aplicado después de un spike;
\item
  \texttt{d} es el incremento de recuperación después del spike.
\end{itemize}

Cuando \texttt{v} alcanza 30 mV se registra un spike y se aplica:

\$ v=c,\qquad u=u+d \$

La elección de este modelo permite representar diferentes patrones de
disparo con pocas variables y un costo computacional reducido.

\hypertarget{generaciuxf3n-de-inputs}{%
\section{6. Generación de inputs}\label{generaciuxf3n-de-inputs}}

Se implementaron distintas formas de corriente externa para estudiar
cómo cambia la actividad neuronal frente a comandos temporales
diferentes:

\begin{itemize}
\tightlist
\item
  corriente constante
\item
  pulso
\item
  rampa
\item
  señal sinusoidal
\item
  comando motor artificial con subida, mantenimiento y relajación
\item
  ruido individual
\end{itemize}

El input se genera como un vector alineado con el tiempo de simulación.
En los pools neuronales puede compartirse una señal común y agregarse
ruido independiente a cada neurona.

\hypertarget{pool-neuronal-y-heterogeneidad}{%
\section{7. Pool neuronal y
heterogeneidad}\label{pool-neuronal-y-heterogeneidad}}

El modelo individual se extendió a poblaciones pequeñas. Cada neurona
conserva sus propias variables \texttt{v} y \texttt{u}, por lo que las
matrices de estado tienen la forma:

\begin{lstlisting}
(número de neuronas, número de pasos temporales)
\end{lstlisting}

La heterogeneidad se introduce mediante pequeñas variaciones aleatorias
en \texttt{a}, \texttt{b}, \texttt{c} y \texttt{d}. Esto evita que todas
las neuronas respondan de manera idéntica al mismo comando motor.

\hypertarget{semilla-aleatoria}{%
\subsection{7.1. Semilla aleatoria}\label{semilla-aleatoria}}

La semilla controla el generador pseudoaleatorio empleado para producir:

\begin{itemize}
\tightlist
\item
  variaciones de parámetros;
\item
  ruido de input;
\item
  máscaras de conectividad.
\end{itemize}

No es una propiedad biológica. Su función es metodológica:

\begin{lstlisting}
misma configuración + misma semilla → misma realización

misma configuración + otra semilla  → nueva realización aleatoria
\end{lstlisting}

Fijar la semilla permite reproducir una simulación y comparar escenarios
sin cambiar simultáneamente el ruido o la conectividad. Repetir un
experimento con varias semillas permite evaluar la robustez del
procedimiento.

\hypertarget{conectividad-sinuxe1ptica}{%
\section{8. Conectividad sináptica}\label{conectividad-sinuxe1ptica}}

\hypertarget{matriz-de-pesos}{%
\subsection{8.1. Matriz de pesos}\label{matriz-de-pesos}}

Las conexiones entre neuronas se representan mediante una matriz
\texttt{W} con la convención:

\begin{lstlisting}
W[i, j] = efecto de la neurona origen i sobre una neurona destino j
\end{lstlisting}

Por lo tanto:

\begin{itemize}
\tightlist
\item
  las filas representan neuronas de origen;
\item
  las columnas representan neuronas de destino;
\item
  un cero indica ausencia de conexión;
\item
  la magnitud indica la fuerza de la conexión;
\item
  el uso o signo del peso determina si la conexión es excitatoria o
  inhibitoria.
\end{itemize}

Las matrices se crean mediante máscaras aleatorias controladas por una
probabilidad de conexión. La probabilidad define cuántas conexiones
existen, mientras que el peso define la magnitud de cada conexión.

\hypertarget{circuito-de-inhibiciuxf3n-recurrente-tipo-renshaw}{%
\section{9. Circuito de inhibición recurrente tipo
Renshaw}\label{circuito-de-inhibiciuxf3n-recurrente-tipo-renshaw}}

El circuito utiliza dos poblaciones:

\begin{itemize}
\tightlist
\item
  un pool de motoneuronas
\item
  un pool más pequeño de células de Renshaw.
\end{itemize}

Las motoneuronas utilizan inicialmente un preset
\texttt{Regular\ Spiking}. Las Renshaw utilizan \texttt{Fast\ Spiking}
como aproximación funcional a una interneurona rápida. Estos presets no
constituyen una calibración biológica específica.

\hypertarget{flujo-del-circuito}{%
\subsection{9.1. Flujo del circuito}\label{flujo-del-circuito}}

\begin{lstlisting}
comando motor
      ↓

pool de motoneuronas

      ↓ 

spikes y estados sinápticos W_MN_to_R

      ↓ 

excitación pool de células de Renshaw

      ↓ 

spikes y estados sinápticos W_R_to_MN

      ↓ 

corriente inhibitoria resta sobre el input de las motoneuronas

      ↓
nueva actualización del pool motor
\end{lstlisting}

\hypertarget{matrices-del-circuito}{%
\subsection{9.2. Matrices del circuito}\label{matrices-del-circuito}}

Se emplean dos matrices distintas:

\begin{lstlisting}
W_MN_to_R:  motoneuronas → Renshaw
W_R_to_MN:  Renshaw → motoneuronas
\end{lstlisting}

Por ejemplo con 20 motoneuronas y 5 Renshaw:

\begin{lstlisting}
W_MN_to_R.shape = (20, 5)
W_R_to_MN.shape = (5, 20)
\end{lstlisting}

Ambas matrices almacenan magnitudes positivas. \texttt{W\_MN\_to\_R} se
utiliza como excitación. \texttt{W\_R\_to\_MN} representa una vía
inhibitoria porque su corriente se resta explícitamente del input motor:

\$ I\_\{MN\}=I\_\{motor\}+I\_\{ruido\}-I\_\{Renshaw\} \$

\$ I\_R=I\_\{MN\rightarrow R\}+I\_\{ruido,R\} \$

Esta convención mantiene separado el valor almacenado en la matriz del
efecto funcional que produce en la ecuación.

\hypertarget{diseuxf1o-de-comparaciones}{%
\section{10. Diseño de comparaciones}\label{diseuxf1o-de-comparaciones}}

Para estudiar el papel de la inhibición recurrente se definieron
escenarios con distintas magnitudes de la vía
\texttt{Renshaw\ →\ motoneurona}.

La metodología mantiene constantes:

\begin{itemize}
\tightlist
\item
  tamaño de los pools
\item
  input externo
\item
  parámetros neuronales base
\item
  nivel de ruido
\item
  probabilidades de conexión
\end{itemize}

El único parámetro que se modifica entre escenarios es la magnitud
inhibitoria.

\hypertarget{muxe9tricas-implementadas}{%
\section{11. Métricas implementadas}\label{muxe9tricas-implementadas}}

Las métricas se utilizan para resumir la actividad sin depender
únicamente de una inspección visual:

\begin{itemize}
\tightlist
\item
  spikes totales
\item
  firing rate medio
\item
  cantidad de motoneuronas activas
\item
  actividad poblacional por ventanas temporales
\item
  pico de actividad poblacional
\item
  corriente inhibitoria media
\item
  intervalos entre spikes
\item
  variabilidad de los intervalos
\item
  fuerza máxima
\item
  fuerza media
\item
  tiempo hasta el pico de fuerza
\end{itemize}

Estas métricas forman parte del método de análisis. Deben interpretarse
junto con rasters, potenciales y señales temporales, ya que cada una
resume un aspecto diferente de la simulación.

\hypertarget{conversiuxf3n-de-spikes-en-fuerza-muscular}{%
\section{12. Conversión de spikes en fuerza
muscular}\label{conversiuxf3n-de-spikes-en-fuerza-muscular}}

\hypertarget{twitch-elemental}{%
\subsection{12.1. Twitch elemental}\label{twitch-elemental}}

Cada spike genera una respuesta contráctil elemental:

\$
h(t)=A\left(e\^{}\{-t/\tau\emph{\{decay\}\}-e\^{}\{-t/\tau}\{rise\}\}\right),\qquad t\geq0
\$

Se utiliza esta diferencia de exponenciales porque produce una subida
rápida seguida por una relajación más lenta. La condición
\texttt{tau\_decay\ \textgreater{}\ tau\_rise} mantiene una forma
positiva y causal.

Se eligió por su simplicidad, bajo número de parámetros y capacidad para
representar sumación temporal sin introducir un modelo muscular
complejo.

\hypertarget{convoluciuxf3n-de-spikes-y-twitches}{%
\subsection{12.2. Convolución de spikes y
twitches}\label{convoluciuxf3n-de-spikes-y-twitches}}

La matriz de spikes se suma entre motoneuronas para formar un tren
poblacional de impulsos. Luego se convoluciona con el kernel twitch:

\$ F{[}t{]}=\sum\_\{k=0\}\^{}\{t\} impulsos{[}k{]}h{[}t-k{]} \$

El procedimiento completo es:

\begin{lstlisting}
Matriz de spikes

      ↓ 
Sumar motoneuronas tren poblacional de impulsos

      ↓ 
Aplicar un twitch a cada impulso convolución temporal

      ↓
Suma de twitches superpuestos

      ↓
Fuerza total en unidades arbitrarias
\end{lstlisting}

La señal obtenida es un proxy de activación muscular. No se calcula
fuerza física, torque articular ni movimiento.


\end{document}
